\documentclass[12pt, a4paper]{ctexart}

\usepackage{geometry}
\geometry{left=2.5cm, right=2.5cm, top=3cm, bottom=3cm}

\usepackage{setspace}
\onehalfspacing

\usepackage{parskip}
\setlength{\parindent}{2em}

\usepackage{titlesec}
\titleformat{\section}{\large\bfseries\centering}{}{0em}{}
\titlespacing{\section}{0pt}{2.5em}{1.5em}

\usepackage{fancyhdr}
\pagestyle{fancy}
\fancyhf{}
\fancyfoot[C]{\thepage}
\renewcommand{\headrulewidth}{0pt}

\title{\textbf{我想活出怎样的人生}}
\author{}
\date{}

\begin{document}

\maketitle
\thispagestyle{empty}

\section*{一}

我出生在南方一座不大不小的城市里，认识我的人都叫我小梅。

我小时候长得很可爱哦，是我发自内心觉得自己可爱。爸爸是商人，妈妈是老师，还有我最喜欢的姐姐，姐姐大我四岁，姐姐很漂亮呢，有很多人喜欢姐姐，我也非常喜欢姐姐。

小学的时候，一切都是那么幸福，那时候感觉姐姐什么都会呀，无论是什么，姐姐都能够教给我。无论是数学题，钢琴，还是乒乓球。姐姐很喜欢教我弹钢琴，她说我的手很柔软很细长，很适合弹钢琴。

那时候的我，觉得自己会一直这样长大。

\section*{二}

我最好的朋友是住在隔壁的竹马君，竹马君长得很帅气，他的性格淡淡的，不像别的男孩子那么淘气，可是竹马君呀，要我说真是一个不开窍的笨蛋。我换了新的发卡问他有没有注意到我有什么不一样的时候，他想了半天问我是不是长高了；他还总是拉着我去打球，真是的哟，竹马君，明明我在意的不是球。

我说不清什么时候开始喜欢他的。那种喜欢很干净，像刚洗过的白衬衫，带着阳光的味道。

\section*{三}

初中，我考上了市里最好的中学。课业压力像潮水一样涌来，排名、考试、竞争，压得我喘不过气。我开始像很多人一样失眠，开始对着卷子发呆，开始在别人问我“怎么了”的时候说“没事”。

我在等别人找到我。

姐姐总是能找到我，她把我当成小猫一样搂在怀里，说我好可爱，告诉我学习并不是唯一的出路，她教我用钢琴弹我喜欢的歌，她教我唱我喜欢的歌，她教会了我那些我看不懂的数学符号背后的含义。

竹马君还是一个笨蛋。他每天放学都陪我走，给我讲愚蠢的冷笑话，在我情绪最低落的那天，硬拽着我去打篮球。我不会打，他就把球塞到我手里，说“随便扔，反正我也接不住”。我扔了几个，一个都没进，却不知道为什么笑出了声。

真是个可爱的孩子，想必会有可爱的生活吧。

\section*{四}

高中，我以为自己终于熬过了所有阴霾。

姐姐考上了大学，我和姐姐没法再天天见面，但姐姐会发消息问“最近怎么样”，偶尔寄回来一些漂亮的东西。竹马君和我上了同一所高中，但还是不开窍的样子。

一切都在变好。

然后，那通电话来了。

那天是周六，我正在学校自习。老师把我叫出去，老师的脸色阴沉。我不记得自己是怎么去的医院，走廊里有刺鼻的消毒水味道，亲戚们都在一旁七嘴八舌。医生走出来，摇了摇头。

我的人生大概停在了这一刻，我大概用了余生的长度来听懂医生在说什么。

假的吧，我蹲了下来，我没有哭，因为我的大脑死机了。

爸爸妈妈和姐姐，在旅游的路上出了车祸。三个人，一个都没留下。

那一年，我十七岁。

\section*{五}

葬礼那天下了很大的雨。我穿着黑色的衣服，站在墓碑前，觉得自己像一具空壳。

我把自己关在家里。亲戚们轮流来敲门。邻居阿姨送饭来。我都没有开门，我抱着姐姐的衣服缩在角落，从早坐到晚，从天黑坐到天亮。我不再说话，不再出门，不再回任何人的消息。

我不想活了。这个念头像野草一样疯长，占据了整个脑子。

只有一个人没有被推开——竹马君。我不知道为什么，只有他的消息我会看，只有他说的“今天食堂的菜好难吃”“你姐以前借我的那本书我还没还”能让我觉得，外面好像还有一点声音。

他每天都来。敲敲门，但他知道我不会开的，他就坐在门外，自言自语。有时候说学校的事，有时候说他们小时候的事。有时候什么都不说，只发一张晚霞的照片。

我坐在门的另一边，抱着膝盖，把耳朵贴在门板上，静静地听。

\section*{六}

我开始重新学说话。

声音很轻，很害怕，像刚学飞的鸟。我隔着一扇门，从嗯、啊，到简单的句子。“然后呢？”“真的吗？”再后来，我问了一句“你今天没上课吗？”

门外沉默了几秒。然后我听见竹马君用压不住颤抖的声音说：“请假了。”

我没有开门，但我会笑了。

然后有一天，竹马君说了很多很多，说到很晚也没走。他说完了所有能说的话，忽然安静了很久。然后，他用一种很轻、很闷的声音说：

“我要搬走了。”

我愣住了。世界好像在一瞬间崩塌。

就在我觉得自己又要掉回那个深渊的时候，竹马君忽然开始说话，语无伦次地说了很多很多喜欢我的话。他说他从小学就开始喜欢我，说他知道现在说这些很不合时宜，说他走了也会回来看我，说他不是要逼我。

我大声哭了出来。哭到整个人都在发抖。

然后，我打开了门。

我冲了出去，用尽全身力气抱住了他。

\section*{七}

我们在一起了。

但我已经落下了太多功课，只能留一级。竹马君去了别的城市读大学，异地恋开始了。

一开始，一切都很好。竹马君每天发消息，视频电话从不间断。他笨拙地安慰我，说异地没关系，说只要我还喜欢他就行。可我看着他的学校里那些优秀、漂亮、自由的女孩子，焦虑像藤蔓一样缠住了我的心。

我拼命地学，我只想考到他身边去。

但是对失败的恐惧始终牵动着我的心弦，我的状态开始变得越来越差，我没有办法长时间地冷静下来做同一件事情。

奇迹没有发生。或许我已经连参加考试的勇气都没有了。

我想了很久，我觉得我是他的累赘，和我在一起他是得不到幸福的。

我割腕了，但是我还是太怕死了，我没死成。他立刻赶回来看我，他向我道歉了很多，说他对我关心不够，他还说他已经退掉了现在正在参加的比赛，全心全意地照顾我。

我看着憔悴的他，觉得我果然是一个累赘。

我和他说了分手，当然以他的个性当然不会接受，所以我删掉了他的所有联系方式，他来找我的所有尝试都被我躲开，因为我知道，男人是聪明的动物，只要我骗他一次，他就永远不会被骗了，我决定离开他，离开这里，那天晚上，我好像流光了这辈子剩下的所有眼泪。

\section*{八}

我决定用余下的每一天惩罚自己。

我是累赘，是不幸的源头，是一切美好的破坏者。我不配拥有任何美好的东西。

从哪里开始惩罚好呢，我看向了镜子里那张依旧可爱的脸庞，我开始出卖自己的身体。

我像玩具一样可怜兮兮的可爱样子似乎很得某些有钱人喜欢，我成为了某个有钱人的专属玩物，我叫他主人，他不光自己玩我，还把我分享给他的朋友们，但是他给了我很多钱，我染上了脏病，没关系，可以去医院治好。怀孕了就去打胎。主人把我挂上了成人视频网站，我很有人气哦，果然我还是很可爱吧？

我开始忘记我是谁，忘记时间的流逝，忘记一切。

\section*{九}

三十岁那年，我被主人抛弃了。新的玩具更年轻，更有新鲜感。我被扔在了街上，口袋里只剩一些钱。

我决定去死。

我沿着河走，想找一个安静的地方，能让我在还算可爱的时候离开。我觉得自己活得太久了，太脏了，太累了。

然后，一只小猫出现在我脚边。

小猫很可爱，毛茸茸的，用圆溜溜的眼睛看着我。我蹲下来，摸了摸它的头。小猫叫了一声，转身往前走。我不知道为什么会跟着它走，也许是因为我太久没有摸过温暖的东西了。

我跟着小猫走了很久，走进了一片别墅区，钻进了一栋房子的院子里。

然后，门开了。

走出来的人，是竹马君。

\section*{十}

他三十岁了，是一个成熟的男人了，看到他的那一刻，我的心都要碎了，但我很高兴他现在过着这么成功和幸福的生活，我想着我现在出现只是在破坏他的家庭，但是已经来不及了，他已经看向了我，我向他打了个招呼，转身就想走，我转身的那一刻，久违的眼泪又一次流了下来。

但是他飞奔着跑过来，抱住了我。

他哭着喊我的名字，一遍又一遍。

我浑身僵硬，脑子一片空白，他请我进入他的家里，就这样，这十年以来我第一次有尊严地坐到了富人区别墅的餐桌上，而不是作为被享用的餐品。

竹马君说了很多，他说他这十年来如何奋斗，如何如何寻找我，讲到动情处，他的眼泪流了下来。

然而我唯一的想法就是，眼前的这个男人不能再和我有一点瓜葛。但我听到他的讲述，我不禁流泪，我过了很久才能够说话，我说你别再想我，我把我这十年以来做的肮脏龌龊的事情都说了，我说的时候被自己逗笑了，我想我和网上说的青楼文学有什么区别，我不能再做一个求大哥带我走的xx了，于是我求他绝情，求他不要再理我这样一个下贱的女人。理我这样一个下贱的女人，对你有什么好处，你会被说闲话，你的生意说不定也会受到影响，你会被网上那些人嘲笑一辈子，你要是和我结婚，你的孩子也一定会被霸凌。

竹马君没有露出我想象中的震惊表情，相反他的语气更坚定了。

他说他找了十年，说他知道一些事，说他不在乎那些。

然后，他说：“小梅，我们一起逃走吧。”

“我们一起，逃走吧。”

这句话让我的大脑停了下来。

“去没有人认识我们的异国他乡，去没有人会议论我们的地方，去你以前一直提起的海滨小镇，去乡下也好。我会说很多语言，我总能想到赚钱的方法。我会好好保护你。”

我哭成了一个泪人。

\section*{十一之一}

第二天早上，竹马君走进我的房间。

我相信我们的余生或许真的可能会是幸福的，如果他走进房间看到的不是我的尸体的话。

但是已经来不及了，他会看到我留给他的那一封长长的信，里面充满了我的忏悔，从小到大我想对他说的每一句话都在里面了。不过果然无论怎样都于事无补吧，我无法想象他看到我尸体的时候会怎样，我不敢想象。

对不起，竹马君，可是我已经没有勇气活下去了，对不起，我不是有意想把你也拉下去。

对不起，对不起，对不起。

我想象得到，因为如果换过来，我是竹马君的话，我的人生在那一天也就结束了，我的余生每一天都已经变成了对那一天的复盘，我会止不住地想如果那一天，如果，如果我那天晚上和小梅一起睡觉，或者就算是不睡觉只是看着她，会不会一切都不一样了呢？

但是已经没有如果了。

对不起，竹马君，我把一切都搞砸了。

\section*{十一之二}

（作者写到这里实在是蚌埠住了，有点破防，于是补充了这个结局，大家可以挑选自己喜欢的作为你认为的结局）

第二天早上，竹马君走进我的房间。

我还在呼吸。阳光透过窗帘照进来，小猫蜷在床脚。我坐在床上，眼睛红肿着。

我看着他，说了一句很轻很轻的话：“我饿了。”

竹马君愣了一下，然后笑了。他笑得比小时候更笨，但更温柔。他说：“我去给你做早餐。”

他们在某个海滨小镇租了一间带院子的房子。不大，但能看到海。竹马君确实学会了当地的语言，在镇上找了一份工作，还在院子里种了番茄。我花了很长时间才敢出门，花了更长时间才不再在镜子前发抖。

我还是做噩梦，还是会在深夜忽然觉得喘不过气。但每次醒来，身边都有一个人。那个笨拙的、十年都没有放弃的人，会握住我的手，轻声说“我一直在这里”。

有一天傍晚，我们坐在海边看夕阳。海风把我的头发吹乱了。

我问：“喂，笨蛋，你后悔吗？”

竹马君转过头，还是那个虎牙，还是那个傻笑。

“我不后悔，我觉得我赚翻了。”

我没有说话。我靠在他的肩上，闭着眼，听着海浪的声音。我闭上眼睛，好像回到了我与他重逢那天的那个晚上，我多么庆幸在那个晚上我没有结束自己的生命，否则此刻一定看不见这么美的夕阳，最重要的是，一定看不到，这么帅气的竹马君，和那个在他的眼中，可爱的，我。

我爱着。我被爱着。

夕阳把海面染成金色，我们是两个人。

\begin{center}
（全文完）
\end{center}

\end{document}